\section{¿Las matemáticas se descubren o se crean?}

Al inicio de la entrevista se discute «lo creado y lo descubierto». Este es el trasfondo filosófico de AI for Math. Si las matemáticas se descubren, la IA necesita aprender estructuras objetivas: las relaciones estables entre números primos, simetrías, topología, curvatura y demostraciones. Si las matemáticas se crean, la IA necesita aprender cómo los seres humanos eligen definiciones, axiomas, notaciones y lenguajes teóricos. En la realidad ambas cosas se entrelazan: los objetos matemáticos tienen restricciones objetivas, mientras que el lenguaje matemático y los caminos de demostración los crean las personas.

\begin{figure}[H]
\centering
\includegraphics[width=0.92\textwidth]{math-created-discovered.png}
\caption{Dos perspectivas de las matemáticas: la estructura descubierta y el lenguaje creado.}
\end{figure}

\begin{knowledgebox}{Lectura de la figura: por qué AI for Math enfrenta dos niveles de tarea a la vez}
El lado izquierdo destaca la objetividad de la estructura matemática; el lado derecho, el carácter creativo del lenguaje matemático y de los sistemas axiomáticos. AI for Math no puede limitarse a la coincidencia de patrones ni a generar texto elegante; debe captar estructuras estables y, al mismo tiempo, escribir el razonamiento en un lenguaje verificable.
\end{knowledgebox}

\subsection{bounded attention y free attention}

En la entrevista aparece la expresión bounded vs free attention, que puede entenderse como dos modos del pensamiento matemático: la bounded attention avanza localmente en torno a definiciones, condiciones, objetivos y herramientas existentes; la free attention, en cambio, permite la asociación, la estética, la analogía y los saltos. La investigación matemática necesita alternar entre ambas: la demostración formal se inclina hacia lo bounded, y la generación de conjeturas e intuiciones, hacia lo free.

\begin{figure}[H]
\centering
\includegraphics[width=0.92\textwidth]{bounded-free-attention.png}
\caption{La bounded attention y la free attention en la intuición matemática.}
\end{figure}

\begin{knowledgebox}{Lectura de la figura: a qué sirve cada tipo de atención}
La bounded attention es adecuada para la revisión de demostraciones, el razonamiento local y la formalización en Lean; la free attention es adecuada para descubrir patrones, plantear conjeturas y establecer analogías entre distintos campos. Un sistema sólido de AI for Math necesita que ambas colaboren, en lugar de saber solo hacer búsquedas rigurosas o solo hacer asociaciones divergentes.
\end{knowledgebox}

\begin{warningbox}{No hay que contraponer la intuición y la formalización}
La intuición matemática no es enemiga de la formalización. Un buen sistema de formalización convierte la intuición en objetos verificables; una buena intuición, a su vez, ayuda a elegir los problemas y las rutas de demostración que vale la pena formalizar.
\end{warningbox}

\subsection{Resumen de la sección}

La dificultad de AI for Math proviene de una doble estructura: las matemáticas se parecen tanto a un mundo objetivo como a una obra de ingeniería del lenguaje humano. El modelo debe alternar continuamente entre la asociación libre y la verificación rigurosa.
